% Generated draft: requires mathematical and source review.
\documentclass[11pt]{article}
\usepackage[margin=25mm]{geometry}
\usepackage{fontspec}
\setmainfont{Latin Modern Roman}
\usepackage{amsmath,amssymb,mathtools,unicode-math}
\setmathfont{Latin Modern Math}
\usepackage{enumitem,needspace,xurl,hyperref}
\setlength{\parindent}{0pt}
\setlength{\parskip}{5pt}
\setlength{\emergencystretch}{3em}
\newcommand{\N}{\mathbb{N}}
\newcommand{\Z}{\mathbb{Z}}
\newcommand{\Q}{\mathbb{Q}}
\newcommand{\R}{\mathbb{R}}
\newcommand{\C}{\mathbb{C}}
\newcommand{\F}{\mathbb{F}}
\newcommand{\A}{\mathbb{A}}
\newcommand{\PP}{\mathbb{P}}
\newcommand{\E}{\mathbb{E}}
\newcommand{\eps}{\varepsilon}
\newcommand{\dd}{\,\mathrm{d}}
\begin{document}
\section*{UnsolvedMath ID: 1100310}
\textbf{Source problem code:} catalogue entry
\subsection{Definitions and mathematical statement}
Let \(F_n\) be the free group of finite rank \(n\geq2\), considered in the pure first-order language of groups. A subset \(H\subseteq F_n\) is definable with parameters if there are a group-language formula \(\varphi(x;y_1,\ldots,y_m)\), \(m\geq0\), and parameters \(a_1,\ldots,a_m\in F_n\) such that
\[
H=\{g\in F_n:F_n\models\varphi(g;a_1,\ldots,a_m)\}.
\]
The conjecture asks whether every such \(H\) that is a subgroup satisfies
\[
H=F_n\quad\text{or}\quad \exists h\in F_n:\ H=\langle h\rangle.
\]
The trivial subgroup is included as cyclic by taking \(h=1\). No parameter-free, finite-index or finite-generation restriction on \(H\) is imposed. The rank-one case is already cyclic and adds no nonabelian case.

The supplied nested research summary reports an affirmative theorem, first obtained from Kharlampovich--Myasnikov's definable-set results (November 2011 preprint; 2013 publication) and independently reproved by Perin--Pillay--Sklinos--Tent, Theorem 4.3 (2014), in the stronger torsion-free hyperbolic setting. It connects the result to Malcev's 1965 question. These are supplied literature claims, not a new proof.

The registry labels the item solved and SOLVED-IN-LITERATURE. The older source was manually transcribed from the 2004 list, accessed 2026-07-29, and its original status/rights evidence remains NEEDS\_REVIEW. This draft preserves those provenance qualifications and does not claim an independent source or current-status check.
\subsection{Short English statement}
Every first-order definable subgroup of a nonabelian finite-rank free group, allowing parameters, should be either cyclic or the whole group. The supplied later literature summary reports this theorem proved, including a stronger torsion-free hyperbolic version.
\subsection{Sources}
\begin{itemize}
\item Local UnsolvedMath registry, permanent ID 1100310, source code catalogue entry, ``Questions in Geometric Group Theory --- Q 3.10,'' including its nested research summary and release-review notes.
\item Bestvina, \emph{Questions in Geometric Group Theory}, listed as 2004, Question 3.10, PDF p. 11; access date 2026-07-29: \url{https://www.math.utah.edu/~bestvina/eprints/questions-updated.pdf}.
\item Kharlampovich--Myasnikov, supplied definable-set reference: November 2011 arXiv preprint and \emph{International Journal of Algebra and Computation} 23 (2013), 91--110. The source supplies no exact title or preprint identifier.
\item Perin--Pillay--Sklinos--Tent, \emph{M\"unster Journal of Mathematics} 7 (2014), Theorem 4.3, as cited in the supplied summary; no title or full page range is supplied.
\item Malcev's 1965 question is mentioned historically in the supplied summary without a full bibliographic reference.
\end{itemize}
\end{document}
